\chapter{Có và Không}
\markboth{Có và Không}{Having and Not Having}

\section{Một người không nhiều tiền nhưng có giờ ăn cơm cùng gia đình.}
\begin{truyen}{Một người không nhiều tiền nhưng có giờ ăn cơm cùng gia đình.}{Not having much can still mean having enough.}
\chuhoa{A}{nh từng ngại khi thấy bạn bè mua xe mới, đổi điện thoại mới. Nhưng tối nào anh cũng có mặt bên mâm cơm, nghe con kể chuyện lớp học, nghe mẹ hỏi vài câu nhỏ. Có những thứ không cầm được trong tay mà lại chính là phần làm nên cảm giác đủ đầy.}
\textit{(He used to feel small when friends bought new cars and new phones. Yet every evening he was present at the family table, hearing his child talk about school and his mother asking simple questions. Some things cannot be held in the hand, yet they are exactly what make life feel full.)}
\end{truyen}
\begin{baihoc}
Thiếu vật chất không đồng nghĩa với thiếu giá trị sống.
\textit{(Material lack does not always mean a lack of life's value.)}
\end{baihoc}
\begin{ghinhoanh}
\textbf{Short English to remember:}

Enough is sometimes invisible.
\end{ghinhoanh}
\begin{tuhocnhanh}
1. Điều gì em đang có nhưng lâu nay xem như bình thường? \textit{(What do you already have that you have started treating as ordinary?)}

2. Nếu bớt so sánh một chút, em thấy đời mình còn thiếu nhiều đến vậy không? \textit{(If you compared less, would your life still feel that lacking?)}
\end{tuhocnhanh}
\ngancach

\section{Một học sinh có đủ sách vở nhưng không còn hứng thú học.}
\begin{truyen}{Một học sinh có đủ sách vở nhưng không còn hứng thú học.}{Having tools is not the same as having fire.}
\chuhoa{E}{m có bàn học đẹp, tài liệu dày, khóa học không thiếu. Nhưng mỗi lần mở sách, em chỉ thấy nặng. Đến lúc đó em mới hiểu điều mình thiếu không phải là dụng cụ, mà là lý do để ngồi xuống học một cách có ý nghĩa.}
\textit{(The student had a good desk, thick materials, and many courses. But every time the books opened, only heaviness appeared. Then came the realization: what was missing was not equipment, but a reason to sit down and learn meaningfully.)}
\end{truyen}
\begin{baihoc}
Có điều kiện học chưa chắc đã có động lực học.
\textit{(Having learning conditions does not guarantee having learning motivation.)}
\end{baihoc}
\begin{ghinhoanh}
\textbf{Short English to remember:}

Tools cannot replace purpose.
\end{ghinhoanh}
\begin{tuhocnhanh}
1. Điều gì đang làm em học mà không còn thấy vui? \textit{(What is making your learning lose its joy?)}

2. Em cần thêm tài liệu hay cần tìm lại lý do học? \textit{(Do you need more material, or do you need to recover your reason for learning?)}
\end{tuhocnhanh}
\ngancach

\section{Một lớp học không có nhiều thiết bị nhưng có tinh thần cùng tiến.}
\begin{truyen}{Một lớp học không có nhiều thiết bị nhưng có tinh thần cùng tiến.}{A room can be poor and still be rich.}
\chuhoa{C}{ăn phòng học cũ, máy chiếu đôi lúc chập chờn, quạt kêu hơi lớn. Nhưng trong lớp ấy, học sinh biết nhường nhau phát biểu, biết hỏi khi chưa hiểu, biết kéo bạn đang nản đứng dậy. Cái lớp thiếu tiện nghi ấy lại có thứ mà nhiều nơi đủ đầy chưa chắc có: khí hậu học tập tử tế.}
\textit{(The classroom was old, the projector unstable, and the fan a little noisy. Yet in that room students gave each other space to speak, asked when they did not understand, and lifted friends who felt discouraged. That under-equipped class had something many well-equipped places still lack: a decent learning climate.)}
\end{truyen}
\begin{baihoc}
Không phải nơi nào ít điều kiện cũng ít hy vọng.
\textit{(Limited conditions do not always mean limited hope.)}
\end{baihoc}
\begin{ghinhoanh}
\textbf{Short English to remember:}

A poor room can hold rich hearts.
\end{ghinhoanh}
\begin{tuhocnhanh}
1. Em đang nhìn lớp mình bằng con mắt thiếu thốn hay bằng con mắt xây dựng? \textit{(Are you looking at your class through the eyes of lack or the eyes of building?)}

2. Một điều nhỏ em có thể góp để lớp tốt hơn là gì? \textit{(What is one small thing you can contribute to make your class better?)}
\end{tuhocnhanh}
\ngancach

\section{Một người không được nhận việc nhưng lại giữ được lòng tự trọng.}
\begin{truyen}{Một người không được nhận việc nhưng lại giữ được lòng tự trọng.}{A closed door can protect your dignity.}
\chuhoa{S}{au buổi phỏng vấn, anh biết mình bị loại. Nhưng anh cũng biết nếu nhận công việc ấy, anh sẽ phải sống bằng cách nịnh bợ và im lặng trước điều sai. Có những cái không đạt được khiến ta buồn trong chốc lát, nhưng giữ cho mình còn nguyên vẹn về lâu dài.}
\textit{(After the interview, he knew he was rejected. Yet he also knew that if he had taken that job, he would have had to live by flattery and silence before what was wrong. Some things we fail to get hurt for a moment, yet protect us for much longer.)}
\end{truyen}
\begin{baihoc}
Không có được một cơ hội chưa chắc là mất; đôi khi đó là được giữ mình.
\textit{(Not obtaining an opportunity is not always loss; sometimes it means keeping yourself.)}
\end{baihoc}
\begin{ghinhoanh}
\textbf{Short English to remember:}

Some losses protect your worth.
\end{ghinhoanh}
\begin{tuhocnhanh}
1. Em có từng tiếc một điều mà sau này thấy may vì không có? \textit{(Have you ever regretted losing something that later proved good not to have?)}

2. Điều gì em không nên đánh đổi để có được thứ mình muốn? \textit{(What should you never trade away in order to get what you want?)}
\end{tuhocnhanh}
\ngancach

\section{Một khoảng trống trong đời khiến người ta lớn lên.}
\begin{truyen}{Một khoảng trống trong đời khiến người ta lớn lên.}{Absence can become a teacher.}
\chuhoa{K}{hi cha đi làm xa, cô bé trong nhà phải tự xếp lại nhiều việc nhỏ. Ban đầu em thấy thiếu, thấy hụt, thấy mọi thứ nặng hơn. Nhưng chính quãng không có đủ đó lại dạy em biết đỡ đần, biết thương mẹ, và biết quý những ngày sum họp hơn trước.}
\textit{(When her father worked far away, the girl had to rearrange many small tasks by herself. At first she felt lack, emptiness, and greater weight. Yet that period of not having enough taught her responsibility, compassion for her mother, and a deeper gratitude for reunions.)}
\end{truyen}
\begin{baihoc}
Cái không đôi khi không chỉ là thiếu, mà còn là khoảng trống để phẩm chất xuất hiện.
\textit{(Not having is sometimes not merely lack, but the space where character appears.)}
\end{baihoc}
\begin{ghinhoanh}
\textbf{Short English to remember:}

Absence can shape character.
\end{ghinhoanh}
\begin{tuhocnhanh}
1. Quãng thiếu thốn nào đã làm em mạnh hơn? \textit{(Which season of lack has made you stronger?)}

2. Em đang nhìn cái không như một bất hạnh hay như một bài học? \textit{(Are you looking at not having as a tragedy or as a lesson?)}
\end{tuhocnhanh}
\ngancach
